\section{Introduction}

\subsection{Background and Context}

A Brain-Computer Interface (BCI), also referred to as a Brain-Machine Interface (BMI), is a system that establishes a direct communication pathway between the human brain and external devices. Its purpose is to allow a person to operate a computer or control another system through brain activity, without relying exclusively on conventional physical input. By interpreting neural signals, a BCI offers an alternative means of interaction for people both with and without physical disabilities.

BCI systems are commonly classified by how brain activity is recorded: invasive, partially invasive, and non-invasive. Invasive systems require electrodes to be surgically implanted inside brain tissue, which allows activity to be recorded close to its source and yields high-quality signals, at the cost of surgical risk. Partially invasive systems, such as electrocorticography (ECoG), place electrodes on the surface of the cerebral cortex without penetrating it. Non-invasive systems record brain activity through sensors positioned on the scalp. They are easier to deploy and avoid surgery altogether, but their recordings are more susceptible to noise, signal attenuation, and other acquisition limitations \citep{ramadan2015basics}.

Electroencephalography (EEG) is the most widely used non-invasive technique. It measures the electrical potentials produced by neural activity through electrodes placed on the scalp. Magnetoencephalography (MEG) is also non-invasive but measures the magnetic fields generated by neural activity and requires large, shielded equipment. EEG is the technique adopted in this research because it is portable and acquires neural signals without surgical intervention.

A BCI system typically operates in several stages. Brain activity is first acquired with the chosen recording technology. The signals are then processed to prepare them for interpretation; depending on the approach, this may involve filtering, artifact removal, feature extraction, and classification. Finally, the interpreted output is translated into commands that are transmitted to and executed by an external device, such as a computer, a wheelchair, a robot, or a remote-controlled vehicle.

Several paradigms can be used to generate and interpret these commands. In motor imagery (MI), the user imagines performing a movement without physically executing it. In steady-state visual evoked potential (SSVEP) systems, the user attends to visual stimuli flickering at different frequencies, each of which produces a corresponding pattern of brain activity. A third approach relies on trained Mental Commands, in which the user trains specific mental actions within a BCI software environment. These paradigms differ in how commands are produced and recognized, and they should not be treated as interchangeable.

Consumer-grade EEG headsets have made these paradigms more accessible to researchers. \citet{tarara2025} explored multiclass motor imagery with the Emotiv EPOC X, classifying six mental states across seven participants. They reported modest improvements in performance after a body-awareness training protocol, but classification remained limited by signal variability within and between participants and by the technical constraints of consumer-grade hardware.

EEG has also been used to control physical vehicles. \citet{setiono2018rccar} built an RC-car control system around the Emotiv EPOC+ headset, combining SSVEP stimulation with Common Spatial Pattern (CSP) feature extraction and a Support Vector Machine (SVM) classifier to recognize three commands: forward, left, and right. Recognized commands were sent wirelessly to an ESP8266 Wemos D1 microcontroller that drove the vehicle. With two participants, the system achieved recognition rates of 88\% for left, 91\% for right, and 90\% for forward, and the authors noted that recognition was affected by electrode contact quality.

Other work has addressed asynchronous control, in which the user decides when to issue a command rather than responding to a cue. \citet{zhao2009asynchronous} used motor imagery to drive a car in a three-dimensional virtual reality environment, using the duration of event-related desynchronization and synchronization (ERD/ERS) as an additional control parameter. Participants were able to steer the virtual car without any muscular activity; the vehicle, however, was simulated rather than physical.

The work closest to this dissertation is the mind-controlled wheelchair developed by \citet{stoyell2022wheelchair}. Their system used the Emotiv EPOC X and the Emotiv Cortex API to recognize trained Mental Commands (push, pull, left, and right). A desktop application processed the command outputs and sent them over a User Datagram Protocol (UDP) connection to a Raspberry Pi, which in turn drove a Sabertooth motor controller through a Universal Asynchronous Receiver-Transmitter (UART) link. The authors reported that command recognition was inconsistent and imprecise, and mitigated this by selecting the most frequently detected command within each half-second window to avoid rapid, unintended changes of direction. The Emotiv software also had to run on a desktop computer, adding a software dependency, and an early version of the system accumulated delays of approximately five to ten seconds until the command-handling script was revised. These observations show that the performance of a BCI-controlled device depends not only on command recognition but also on the processing and communication mechanisms that connect its components.

Taken together, these studies demonstrate that EEG can control both virtual and physical devices, while exposing recurring challenges in signal quality, command recognition, communication, and responsiveness. Building on this work, the present dissertation integrates the Emotiv EPOC X with a physical remote-controlled car using trained Mental Commands and systematically compares communication approaches used to deliver commands to the vehicle.

The headset-to-host connection is distinct from the communication approaches compared in this dissertation. EEG processing and Mental Command generation take place in the supported Emotiv host environment; the protocols under investigation carry commands from the host application to the vehicle controller.

\subsection{Motivation}

Remote-controlled vehicles are normally operated through physical interaction with a controller, such as pressing buttons or moving joysticks. A BCI offers an alternative input method in which commands are derived from brain activity instead. Building such a system on consumer-grade hardware raises practical questions that go beyond signal classification: how the headset's output is accessed, how commands are transmitted, and how quickly and dependably they become physical movement.

The studies reviewed above answer these questions only in part, because they differ in device, control paradigm, software architecture, and experimental setting. \citet{setiono2018rccar} controlled a physical RC car, but with the EPOC+ and a custom SSVEP classification pipeline. \citet{tarara2025} used the EPOC X, but to classify mental states rather than to control a vehicle. \citet{stoyell2022wheelchair} combined the EPOC X and trained Mental Commands to drive a physical device, but that device was a wheelchair, and the work reported problems with consistency and delay.

This motivates an investigation of the EPOC X and its trained Mental Commands as the control source for a small remote-controlled vehicle. Rather than proposing a new EEG classification algorithm, the project focuses on integrating the headset's existing command-recognition functionality with a vehicle-control architecture. It compares communication approaches under defined conditions, measuring latency, delivery, stability and resource use, and evaluates how these characteristics relate to physical vehicle response.

\subsection{Problem Statement}

EEG-based BCI systems have been shown to control external devices, but reliable real-time operation remains a practical challenge. Recognizing brain activity is only the first step: the recognized command must then be communicated to the target device, interpreted by its control system, and translated into the appropriate physical action. Errors and delays can arise at each of these stages.

The literature reflects this. The RC car of \citet{setiono2018rccar} depended on a custom signal-processing pipeline and was tested with only two participants; \citet{tarara2025} found that signal variability and hardware constraints limited EPOC X classification performance; and \citet{stoyell2022wheelchair} encountered inconsistent recognition, software dependencies, and communication delays when connecting the EPOC X to a physical wheelchair. Integrating a BCI with a physical device therefore requires attention to the complete control chain, from command recognition to physical execution.

The problem addressed in this dissertation is the lack of a controlled comparison of communication approaches within an Emotiv EPOC X-based brain-controlled vehicle system. The work therefore develops an integrated car-control platform and experimentally evaluates communication performance and its relationship to end-to-end command execution and vehicle response.

\subsection{Research Gap}

Table~\ref{tab:ch1-related-work} summarizes the studies discussed in Section~1.1 along the dimensions most relevant to this dissertation.

\begin{table}[htbp]
  \centering
  \footnotesize
  \caption{Comparison of related EEG-based control studies.}
  \label{tab:ch1-related-work}
  \begin{tabular}{@{}l l p{2.6cm} p{2.4cm}@{}}
    \toprule
    \textbf{Study} & \textbf{Headset} & \textbf{Paradigm} & \textbf{Controlled device} \\
    \midrule
    \citet{setiono2018rccar}        & Emotiv EPOC+ & SSVEP (CSP + SVM)        & Physical RC car \\
    \citet{zhao2009asynchronous}    & --- & Asynchronous MI          & Virtual car \\
    \citet{stoyell2022wheelchair}   & Emotiv EPOC X & Trained Mental Commands & Physical wheelchair \\
    \citet{tarara2025}              & Emotiv EPOC X & Multiclass MI           & None (classification only) \\
    \midrule
    This dissertation               & Emotiv EPOC X & Trained Mental Commands & Physical RC car \\
    \bottomrule
  \end{tabular}
\end{table}

No reviewed study combines all three elements of the present work. \citet{setiono2018rccar} controlled a physical RC car, but with the EPOC+ and an SSVEP-based CSP and SVM pipeline rather than the EPOC X and its trained Mental Commands. \citet{tarara2025} used the EPOC X but stopped at classification, without connecting the recognized commands to a vehicle. \citet{zhao2009asynchronous} showed that asynchronous motor imagery can steer a car without continuous visual stimulation, but only in a virtual environment.

\citet{stoyell2022wheelchair} come closest, demonstrating a complete architecture from the Emotiv Cortex API through a desktop application, UDP communication, and a Raspberry Pi to a motor controller. Their work, however, targeted wheelchair navigation, reported limitations in command consistency, software dependencies, and responsiveness, and did not provide a systematic quantitative evaluation of recognition accuracy and end-to-end performance across a broad group of participants.

The opportunity identified is therefore to adapt the EPOC X and its trained Mental Commands to a small physical remote-controlled car and to compare communication approaches in a common system and experimental setup. The evaluation covers communication latency, delivery, ordering, throughput, resource use and recovery under defined network conditions, as well as command execution and end-to-end vehicle response. The contribution sought is a practical, reproducible comparison, not a new EEG classification method.

This gap is provisional. It does not imply that no similar implementations exist, and it will be refined through the broader review in Chapter~2 before any definitive claim about the originality of the work is made.

\subsection{Research Question}

The main research question guiding this dissertation is:

\begin{quote}
\textbf{How do different communication protocols and messaging approaches affect the latency, reliability, and responsiveness of an EEG-based remote-controlled vehicle using the Emotiv EPOC X and Mental Commands?}
\end{quote}

The research question is supported by the following secondary questions:

\begin{itemize}
  \item How can the Emotiv EPOC X and its Mental Commands functionality be integrated with a physical remote-controlled vehicle?
  \item How do UDP, TCP, WebSocket, MQTT, and other selected communication approaches compare in terms of latency, jitter, message delivery, and connection stability?
  \item How do MQTT Quality of Service configurations influence communication performance in the proposed application?
  \item To what extent do communication characteristics influence the execution of vehicle-control commands?
  \item Which evaluated approach provides the most appropriate balance between responsiveness, reliability, and operational stability for the proposed system?
\end{itemize}

\subsection{Main Research Objective}

The main objective of this dissertation is to design, implement, and experimentally evaluate an EEG-based remote-controlled vehicle using the Emotiv EPOC X and Mental Commands, with particular emphasis on comparing communication protocols and messaging approaches used to transmit control commands between the host computer and the vehicle controller.

The research aims to identify the performance characteristics and trade-offs of the selected approaches under controlled experimental conditions and determine their suitability for responsive and reliable BCI-based vehicle control.

\subsection{Specific Objectives}

To achieve the main objective, the dissertation pursues the following specific objectives:

\begin{enumerate}
  \item \textbf{Investigate the theoretical foundations of EEG-based BCI systems:} review relevant research on EEG acquisition, BCI paradigms, consumer-grade EEG devices, Mental Commands, and their applications in robotics and vehicle control.
  \item \textbf{Examine communication technologies for real-time control applications:} investigate UDP, TCP, WebSocket, MQTT, HTTP, CoAP, and other relevant approaches, focusing on latency, reliability, message delivery, and resource consumption.
  \item \textbf{Design and develop the brain-controlled vehicle architecture:} integrate the Emotiv EPOC X, Emotiv processing environment, host application, communication framework, vehicle controller, and motor-control components.
  \item \textbf{Implement a common communication framework:} develop a modular software architecture that transmits equivalent vehicle-control messages using a consistent message structure.
  \item \textbf{Implement and configure selected communication approaches:} evaluate UDP, TCP, WebSocket, and MQTT, including relevant MQTT Quality of Service configurations, while considering additional approaches according to feasibility.
  \item \textbf{Establish a controlled experimental environment:} define the hardware, network configuration, message payloads, transmission rates, test conditions, and repetition procedures required for a fair comparison.
  \item \textbf{Measure communication performance:} evaluate latency, round-trip time, jitter, message delivery, message ordering, throughput, connection stability, and recovery behavior.
  \item \textbf{Evaluate physical vehicle-control performance:} investigate how communication characteristics affect command execution and physical vehicle responsiveness.
  \item \textbf{Evaluate the integrated EEG-based control system:} assess command recognition, missed commands, false activations, and end-to-end response time while distinguishing BCI processing delays from communication and vehicle-response delays.
  \item \textbf{Analyze and compare experimental results:} identify the trade-offs between evaluated approaches and determine which configurations are most suitable for the proposed application.
\end{enumerate}

\subsection{Research Hypotheses}

Based on the research objectives, the following hypotheses are proposed for experimental investigation:

\begin{description}
  \item[H1 - Communication latency:] Different communication approaches will produce measurable differences in command transmission latency under equivalent hardware, network, and message-generation conditions.
  \item[H2 - Communication reliability:] Communication approaches with different delivery mechanisms and reliability guarantees will exhibit measurable differences in message loss, duplication, ordering, and connection behavior.
  \item[H3 - MQTT Quality of Service:] Different MQTT Quality of Service configurations will produce measurable differences in delivery behavior and communication latency, reflecting trade-offs between delivery guarantees and protocol overhead.
  \item[H4 - Network conditions:] Increased network traffic, delay, and packet loss will affect the performance of the evaluated communication approaches, although the magnitude of these effects may differ between approaches.
  \item[H5 - Vehicle responsiveness:] Differences in communication performance will influence the time and consistency with which the physical vehicle executes transmitted control commands.
  \item[H6 - End-to-end performance:] The responsiveness of the complete EEG-based vehicle-control system will depend on Mental Command generation, communication, and vehicle actuation; communication performance alone will not fully explain end-to-end response time.
\end{description}

These hypotheses will be examined through controlled experiments. No particular protocol is assumed in advance to outperform all others; the results will determine whether the hypotheses are supported by the collected evidence.

\subsection{Scope and Limitations}

This dissertation focuses on the development and experimental evaluation of a physical remote-controlled car operated through EEG-based commands generated using the Emotiv EPOC X and Emotiv Mental Commands functionality. The scope includes:

\begin{description}
  \item[EEG acquisition and command generation:] The Emotiv EPOC X and supported Emotiv processing environment will generate trained control commands. The research focuses on integrating and evaluating these outputs rather than replacing Emotiv's proprietary command-processing system.
  \item[Communication architecture:] The comparison focuses on application-to-vehicle communication. EEG-derived commands are received by a host application and transmitted to the vehicle controller; the native headset-to-host connection remains separate from the protocols compared.
  \item[Communication approaches:] The core investigation considers UDP, TCP, WebSocket, and MQTT, including relevant MQTT Quality of Service levels. HTTP and CoAP may be included where feasible and justified. Bluetooth/BLE and UART may be investigated as alternative physical communication links, rather than treated as equivalent network-protocol comparisons.
  \item[Physical vehicle control:] A physical remote-controlled car is constructed or adapted using an appropriate microcontroller, motor driver, and electronic components, with recognized Mental Commands mapped to predefined vehicle actions.
  \item[Performance evaluation:] The research evaluates communication latency, jitter, message delivery and ordering, throughput, connection stability, and recovery, as well as command execution, vehicle responsiveness, and end-to-end system performance.
\end{description}

The following limitations are acknowledged:

\begin{itemize}
  \item The consumer-grade EPOC X's signal quality and command-generation performance may be affected by electrode contact, environmental noise, user variability, and fatigue.
  \item Emotiv Mental Commands depends on Emotiv's software ecosystem; the research is constrained by the functionality and data exposed through its supported interface.
  \item The communication comparison is limited to selected protocols, implementations, hardware, and network configurations; results should not be interpreted as a universal ranking of communication technologies.
  \item MQTT performance may depend on broker configuration, deployment location, and Quality of Service settings.
  \item Communication latency measurements may be affected by clock synchronization, operating-system scheduling, network conditions, and measurement instrumentation.
  \item Vehicle response time may be influenced by motor characteristics, battery condition, mechanical components, and firmware execution.
  \item The number and characteristics of participants, sessions, and trials may limit the generalization of EEG-related findings.
  \item The vehicle is an experimental platform and is not intended for safety-critical transportation or clinical mobility.
\end{itemize}

The research will prioritize controlled and reproducible experiments while documenting technical and methodological constraints that may influence the results. For safety, testing will take place in a controlled environment with an independent emergency-stop mechanism; the system is not intended for public roads.

\subsection{Expected Contributions}

This dissertation is expected to contribute to BCI integration, communication performance evaluation, and embedded vehicle control:

\begin{enumerate}
  \item \textbf{An integrated EEG-based vehicle-control system} combining the Emotiv EPOC X, Mental Commands, a host-side processing application, a communication framework, and a physical remote-controlled car.
  \item \textbf{A modular communication framework} that implements different communication approaches using a common command structure to support consistent comparisons.
  \item \textbf{An experimental comparison of communication approaches} measuring latency, jitter, delivery behavior, ordering, throughput, and connection stability for selected implementations.
  \item \textbf{An evaluation of MQTT delivery configurations} examining how Quality of Service levels influence delivery and performance in a time-sensitive vehicle-control application.
  \item \textbf{An evaluation of communication effects on physical control} investigating how network behavior influences command execution and vehicle responsiveness.
  \item \textbf{A separation of BCI and communication performance} distinguishing command-generation and recognition delays from communication and vehicle-actuation delays in the end-to-end response.
  \item \textbf{Experimental findings and implementation guidance} on the suitability and trade-offs of evaluated approaches for EEG-based remote-control and related embedded applications.
\end{enumerate}

\subsection{Dissertation Structure}

The remainder of this dissertation is organized as follows.

\begin{itemize}
  \item \textbf{Chapter 2 - Brain-Computer Interfaces and Communication Technologies for EEG-Based Control:} reviews BCI and EEG fundamentals, the Emotiv EPOC X and Mental Commands, vehicle-control applications, real-time challenges, communication technologies, and the research gap.
  \item \textbf{Chapter 3 - Development of an Emotiv EPOC X-Based Brain-Controlled Car System:} presents the system requirements, architecture, hardware, Python framework, EEG command processing, communication protocol implementations, vehicle control, safety mechanisms, and integration testing.
  \item \textbf{Chapter 4 - Research Methodology and Experimental Evaluation:} describes the research design, equipment, participants, experimental conditions and procedures, evaluation metrics, data analysis, and considerations for reproducibility and safety.
  \item \textbf{Chapter 5 - Results:} reports the communication benchmarks, network-condition tests, physical vehicle-control results, EEG command performance, and end-to-end response times.
  \item \textbf{Chapter 6 - Discussion:} interprets the findings against the research objectives and hypotheses, examines communication trade-offs and practical implications, and discusses limitations and improvements.
  \item \textbf{Chapter 7 - Conclusion and Future Work:} summarizes the research and its contributions, evaluates the achievement of the objectives, and presents limitations and recommendations for future research.
\end{itemize}
